\chapter[Chapter \thechapter]{}
\lettrine[ante=`]{T}{hen}  you won't marry me?' said Lord~Peter.

\zz
The prisoner shook her head.

\zz
<No. It wouldn't be fair to you. And besides\longdash>

\zz
<Well?>

<I'm frightened of it. One couldn't get away. I'll live with you, if you like, but I won't marry you.>

Her tone was so unutterably dreary that Wimsey could feel no enthusiasm for this handsome offer.

<But that sort of thing doesn't always work,> he expostulated. <Dash it all, you ought to know—forgive my alluding to it and all that—but it's frightfully inconvenient, and one has just as many rows as if one was married.>

<I know that. But you could cut loose any time you wanted to.>

<But I shouldn't want to.>

<Oh, yes, you would. You've got a family and traditions, you know. Caesar's wife and that sort of thing.>

<Blast Caesar's wife! And as for the family traditions—they're on my side, for what they're worth. Anything a Wimsey does is right and heaven help the person who gets in the way. We've even got a damned old family motto about it—<I hold by my Whimsy>—quite right too. I can't say that when I look in the glass I exactly suggest to myself the original Gerald de Wimsey, who bucked about on a cart horse at the Siege of Acre, but I do jolly well intend to do what I like about marrying. Who's to stop me? They can't eat me. They can't even cut me, if it comes to that. Joke, unintentional, officers, for the use of.>

Harriet laughed.

<No, I suppose they can't cut you. You wouldn't have to slink abroad with your impossible wife and live at obscure continental watering-places like people in Victorian novels.>

<Certainly not.>

<People would forget I'd had a lover?>

<My dear child, they're forgetting that kind of thing every day. They're experts at it.>

<And was supposed to have murdered him?>

<And were triumphantly acquitted of having murdered him, however greatly provoked.>

<Well, I won't marry you. If people can forget all that, they can forget we're not married.>

<Oh, yes, \emph{they} could. I couldn't, that's all. We don't seem to be progressing very fast with this conversation. I take it the general idea of living with me does not hopelessly repel you?>

<But this is all so preposterous,> protested the girl. <How can I say what I should or shouldn't do if I were free and certain of—surviving?>

<Why not? I can imagine what I should do even in the most unlikely circumstances, whereas this really is a dead cert, straight from the stables.>

<I can't,> said Harriet, beginning to wilt. <Do please stop asking me. I don't know. I can't think. I can't see beyond the—beyond the—beyond the next few weeks. I only want to get out of this and be left alone.>

<All right,> said Wimsey, <I won't worry you. Not fair. Abusing my privilege and so on. You can't say <Pig> and sweep out, under the circs., so I won't offend again. As a matter of fact I'll sweep out myself, having an appointment—with a manicurist. Nice little girl, but a trifle refained in her vowels. Cheerio!>

\divider

The manicurist, who had been discovered by the help of Chief-Inspector Parker and his sleuths, was a kitten-faced child with an inviting manner and a shrewd eye. She made no bones about accepting her client's invitation to dine, and showed no surprise when he confidentially murmured that he had a little proposition to put before her. She put her plump elbows on the table, cocked her head at a coy angle, and prepared to sell her honour dear.

As the proposition unfolded itself, her manner underwent an alteration that was almost comical. Her eyes lost their round innocence, her very hair seemed to grow less fluffy, and her eyebrows puckered in genuine astonishment.

<Why, of course I could,> she said finally, <but whatever do you want them for? Seems funny to me.>

<Call it just a joke,> said Wimsey.

<No.> Her mouth hardened. <I wouldn't like it. It doesn't make sense, if you see what I mean. What I mean, it sounds a queer sort of joke and that kind of thing might get a girl into trouble. I say, it's not one of those, what do they call 'em?—there was a bit about it in Madame Crystal's column last week, in \workname{Susie's Snippets}—spells, you know, witchcraft—the occult, that sort of thing? I wouldn't like it if it was to do any harm to anybody.>

<I'm not going to make a waxen image, if that's what you mean. Look here, are you the sort of girl who can keep a secret?>

<Oh, I don't talk. I never was one to let my tongue wag around. I'm not like ordinary girls.>

<No, I thought you weren't. That's why I asked you to come out with me. Well, listen, and I'll tell you.>

He leaned forward and talked. The little painted face upturned to his grew so absorbed and so excited that a bosom friend, dining at a table some way off, grew quite peevish with envy, making sure that darling Mabel was being offered a flat in Paris, a Daimler car and a thousand-pound necklace, and quarrelled fatally with her own escort in consequence.

<So you see,> said Wimsey, <it means a lot to me.>

Darling Mabel gave an ecstatic sigh.

<Is that all true? You're not making it up? It's better than any of the talkies.>

<Yes, but you mustn't say one word. You're the only person I've told. You won't give me away to him?>

<Him? He's a stingy pig. Catch me giving him anything. I'm on. I'll do it for you. It'll be a bit difficult, 'cause I'll have to use the scissors, which we don't do as a rule. But I'll manage. You trust me. They won't be big ones, you know. He comes in pretty often, but I'll give you all I get. And I'll fix Fred. He always has Fred. Fred'll do it if I ask him. What'll I do with them when I get them?>

Wimsey drew an envelope from his pocket.

<Sealed up inside this,> he said, impressively, <there are two little pill-boxes. You mustn't take them out till you get the specimens, because they've been carefully prepared so as to be absolutely chemically clean, if you see what I mean. When you're ready, open the envelope, take out the pill-boxes, put the parings into one and the hair into the other, shut them up at once, put them into a clean envelope and post them to this address. Get that?>

<Yes.> She stretched out an eager hand.

<Good girl. And not a word.>

<Not—one—word!> She made a gesture of exaggerated caution.

<When's your birthday?>

<Oh, I don't have one. I never grow up.>

<Right; then I can send you an unbirthday present any day in the year. You'd look nice in mink, I think.>

<Mink, I think,> she mocked him. <Quite a poet, aren't you?>

<You inspire me,> said Wimsey, politely.